Un message venu de l'École nationale des chartes a déclenché ce livre : pouvait-on refaire,
sur les manuscrits mathématiques de Leibniz, ce qui avait été fait sur les notes de
Grothendieck et sur les manuscrits de Gallica ? La réponse courte est oui. La réponse
longue tient en ces pages, parce que la méthode change moins d'un fonds à l'autre que les
sources elles-mêmes.

\emph{Mathématiques outillées}, le premier livre, décrit les outils : le dépôt \texttt{git},
les issues, l'intégration continue, Lean, l'assistant et ses procédures écrites. Il ne
nomme aucun corpus. Celui-ci fait l'inverse. Il nomme chaque fonds sur lequel la méthode
a servi, dit comment on y accède, ce qui a résisté (un certificat expiré, un serveur qui
interdit l'affichage dans une autre page, un pare-feu qui bloque, un droit d'auteur qui
court jusqu'en 2084) et ce qu'on a décidé. On y apprend surtout à lire une source avant
de commencer à la transcrire.

La méthode a tourné sur deux dépôts publics. \texttt{grothendieck-project} part du fonds
Grothendieck de l'université de Montpellier et s'est étendu aux rédactions Bourbaki, aux
carnets de Quillen, au fonds Grothendieck--Illusie de l'IHÉS, aux papiers de la BnF et aux
enregistrements des cours de Buffalo. \texttt{germain-project} applique la même chaîne aux
manuscrits de dix-neuf mathématiciens que Gallica sert en entier, à commencer par Sophie
Germain. Les deux sources de ce second projet sont dans le domaine public, et c'est ce
qui le rend proche du cas de Leibniz.

Les chiffres datent de début novembre 2026 et viennent des fichiers des deux dépôts. Les
constats techniques (en-têtes HTTP, licences des manifestes) ont été faits depuis un
terminal à une date donnée à chaque fois. Ils changent : on les refait avant de
s'appuyer dessus.

\begin{noteblock}
Personne ici n'est juriste. Ce que ce livre dit des droits est la position que les deux
projets ont prise, avec ses raisons, et non un avis juridique. Une ou un archiviste, ou le
service juridique de l'institution qui détient le fonds, a le dernier mot.
\end{noteblock}
